\chapter{Reachability, 2-CNF, and First-Order Validity}
\label{sec:p04-decision-problem-hardness}

Completeness of a decision problem requires its upper bound and a reduction
from every problem in the declared class. The linear SCC algorithm for
2-CNF does not by itself give a logarithmic-space lower bound. Similarly,
enumerability and undecidability do not by themselves give c.e.-completeness.
The reductions here supply the stronger conclusions with explicit inputs.
The counting argument follows Immerman's complementation method
\cite{Immerman1988Complement}; the first-order construction gives an explicit
finite transition theory and its countermodel, in the machine-to-logic
tradition of \cite[\S11]{Turing1937}, with the distinct later correction
\cite[pp.~544--545]{Turing1938DecisionCorrection}. It is a new construction,
not a transcription of that historical encoding.

\section{Logarithmic space and directed reachability}
\label{sec:p04-dh-reachability}

\begin{hbdefinition}{def:p04-dh-nl}{Logarithmic-space languages and reductions}
Use the read-only-input and all-branches-halting convention of
\cref{def:p04-complexity-convention}. Put
\[
 \mathsf L=\mathsf{SPACE}(\log(n+2)),\qquad
 \mathsf{NL}=\mathsf{NSPACE}(\log(n+2)),
 \qquad
 \mathsf{coNL}=\{\{0,1\}^*\setminus A:A\in\mathsf{NL}\}.
\]
Here \(n\) is total input bit length; changing the finitely many smallest
lengths does not change the class. A logarithmic-space many-one translator
has a read-only input, a write-only output, and \(O(\log(n+2))\) work cells.
It halts on every input and produces one polynomial-length output word.
The notation is \(A\leq_m^{\log}B\), as in
\cref{def:p04-succinct-problems}. Completeness below uses this reduction,
not merely a polynomial-time translation.
\end{hbdefinition}

\begin{hbdefinition}{def:p04-dh-graphs}{Two explicit directed-graph languages}
An instance has a unary positive vertex count \(N\), binary designated
indices \(s,t\in\{1,\ldots,N\}\), and the explicit \(N\)-by-\(N\) Boolean
adjacency matrix in row-major order. Field separators and bits are tokens
in a fixed finite alphabet, encoded by fixed-length binary blocks; indices
use their canonical binary strings. Syntactic validity, lengths and index
bounds are logarithmic-space checkable.
\[
\begin{aligned}
 \mathsf{STCON}&=\{\langle G,s,t\rangle:
       \text{valid encoding and a directed path }s\leadsto t\},\\
 \mathsf{NONREACH}&=\{\langle G,s,t\rangle:
       \text{valid encoding and no directed path }s\leadsto t\}.
\end{aligned}
\]
Paths include length zero. Both languages reject malformed inputs,
including a zero vertex count. Thus \(\mathsf{NONREACH}\) is the complement
of \(\mathsf{STCON}\) on valid graph encodings, not their whole binary
input space. Each vertex index needs \(O(\log(N+1))\) bits; an adjacency
query is computed by rescanning the input.
\end{hbdefinition}

\begin{hbtheorem}{thm:p04-dh-stcon}{proved=annex:p04.dh.stcon}{Directed reachability is NL-complete}
Under \cref{def:p04-dh-graphs}, \(\mathsf{STCON}\) is
\(\mathsf{NL}\)-complete under logarithmic-space many-one reductions.
The hardness translator outputs a valid graph on every input, including
the empty word.
\end{hbtheorem}

\begin{hbmethod}{def:p04-dh-counting}{The reachability census}
For a valid graph, put
\[
 R_d=\{v:\text{a path from }s\text{ to }v\text{ has length at most }d\},
 \qquad c_d=|R_d|.
\]
Start with \(c_0=1\). Given the correct \(c_d\), test one candidate \(v\)
for membership in \(R_{d+1}\) by scanning each \(u=1,\ldots,N\) once.
Nondeterministically include or omit \(u\). Inclusion must be accompanied
by a verified path of length at most \(d\) from \(s\) to \(u\).
Count inclusions and set a flag if an included \(u\) is \(v\) or has an
edge to \(v\). Reject this branch unless the final inclusion count is
exactly \(c_d\). Its flag then gives the correct membership answer.
Repeat for every \(v\), counting true flags to obtain \(c_{d+1}\).
Repeat through \(d=N-2\), then enumerate \(R_{N-1}\) with the same
count check to test whether \(t\) is absent. For \(N=1\) there is no
census-update iteration.
\end{hbmethod}

\begin{hbtheorem}{thm:p04-dh-complement}{proved=annex:p04.dh.complement}{Nonreachability, complementation, and hardness}
The census decides \(\mathsf{NONREACH}\) in \(\mathsf{NL}\).
Together with the configuration reduction, it gives
\(\mathsf{NL}=\mathsf{coNL}\), and \(\mathsf{NONREACH}\) is
\(\mathsf{NL}\)-complete under \(\leq_m^{\log}\).
These conclusions use the verified census; they are not consequences of
the SCC running-time bound.
\end{hbtheorem}

\section{The exact 2-CNF classification}
\label{sec:p04-dh-two-cnf}

\begin{hbdefinition}{def:p04-dh-two-cnf}{Explicit 2-CNF input}
The input consists of a unary count \(r\geq0\) and an explicit list of
clauses over variables \(p_1,\ldots,p_r\); literal indices are canonical
binary strings and each clause has zero, one, or two signed literals.
Clause and literal separators have the fixed binary token coding above.
Invalid indices, clauses of greater length and malformed words are no
instances. The language \(\mathsf{2SAT}\) contains exactly valid
satisfiable inputs. An empty clause makes a valid input unsatisfiable;
an empty clause list is satisfiable, including \(r=0\).
Unused declared variables are permitted. Unary declaration size makes
\(r\) at most the input length.

For an ordinary explicit list with arbitrary finite binary atom names,
replace each distinct occurring name by its occurrence-order index.
A logarithmic-space transducer compares two names by rescanning their
input positions and finds a name's index by counting first occurrences.
It need not store the names simultaneously. Unused declarations can be
discarded. Thus this explicit encoding is logspace intertranslatable
with the ordinary occurring-variable clause-list presentation, but not
with a succinctly described exponentially large formula.
\end{hbdefinition}

\begin{hbtheorem}{thm:p04-dh-two-sat}{proved=annex:p04.dh.two-sat}{2-CNF satisfiability is NL-complete}
The language \(\mathsf{2SAT}\) is \(\mathsf{NL}\)-complete under
logarithmic-space many-one reductions. Its unsatisfiability problem on
valid encodings is also \(\mathsf{NL}\)-complete. This coexists with the
linear literal-reference SCC algorithm of \cref{thm:p02-two-sat};
no equality \(\mathsf L=\mathsf{NL}\), linear bit-time bound, or succinct
input classification is inferred.
\end{hbtheorem}

The hardness translation has one variable for each graph vertex:
\[
 F_{G,s,t}
 =p_s\land\neg p_t\land
    \bigwedge_{(u,v)\in E(G)}(\neg p_u\lor p_v).
\]
It is satisfiable exactly when \(t\) is not reachable from \(s\).
The two unit clauses are contradictory if \(s=t\), matching length-zero
reachability. If there are no edges and \(s\ne t\), setting only the
reachable vertex \(s\) true is a satisfying assignment. Malformed graph
input maps to the fixed unsatisfiable \(p_1\land\neg p_1\), not to a
promise violation.

\section{An effective first-order validity language}
\label{sec:p04-dh-first-order}

\begin{hbdefinition}{def:p04-dh-validity}{Codes and semantic range of validity}
Fix the countable effective signature
\[
 \Sigma_{\mathrm H}
 =\{\mathsf{nil},\mathsf{Reach}\}
   \cup\{\mathsf a_i:i\in\N\}
   \cup\{\mathsf c_j:j\in\N\},
\]
where \(\mathsf{nil}\) is a constant, \(\mathsf{Reach}\) is unary
relational, every \(\mathsf a_i\) is unary functional, and every
\(\mathsf c_j\) is binary functional. Each sentence uses finitely many
symbols. For an index \(i\geq0\), let \(k\) be the length of
\(\operatorname{bin}(i+1)\) and use the self-delimiting field
\(\gamma(i)=1^k0\operatorname{bin}(i+1)\).
Use four-bit tags for the following prefix grammar; juxtaposition means
concatenation of codes:
\[
\begin{array}{c|l@{\qquad}c|l}
0000&\mathsf{nil}&0001&x_i:\ \gamma(i)\\
0010&\mathsf a_i(t):\ \gamma(i)t&0011&\mathsf c_j(t,u):\ \gamma(j)tu\\
0100&\mathsf{Reach}(t):\ t&0101&t=u:\ tu\\
0110&\neg A:\ A&0111&A\land B:\ AB\\
1000&\forall x_i A:\ \gamma(i)A&&
\end{array}
\]
In this grammar each displayed expression after the colon is the tag's
remaining code, and \(t,u,A,B\) stand for recursively encoded expressions,
not their printed names. Tags not listed are invalid. The parser checks
term/formula categories, field lengths, arities, exhausted input and bound
variable indices. Thus being a closed sentence is decidable.
Displayed implication, disjunction and existential quantification expand
to the primitive \(\neg,\land,\forall\); truth and falsity expand to a
fixed tautology or contradiction using \(\mathsf{Reach}(\mathsf{nil})\).
For a binary string
\(u\), use
\[
 \operatorname{code}(u)=2^{|u|}-1+\operatorname{value}_2(u),
 \qquad \operatorname{value}_2(\epsilon)=0.
\]
The inverse writes the canonical binary expansion of the input number
plus one and deletes its leading \(1\). Thus the empty string has code
zero and is an invalid sentence; every natural number has a unique
decoded binary string. This is an explicit computable sentence coding,
not the unrelated machine-program numbering.
For a formula \(A\), \(\operatorname{code}(A)\) means the code of its
declared prefix word, not the code of its printed TeX. For example the
closed atom \(\mathsf{Reach}(\mathsf{nil})\) has prefix word \(01000000\)
and numerical code \(2^8-1+64=319\).

Let \(\mathsf{VALID}_{\mathrm H}\subseteq\N\) consist of the codes of
closed first-order sentences in this signature true in every nonempty
\(\Sigma_{\mathrm H}\)-structure. Malformed codes, open formulas and
arity errors are outside the language. Identity may be included with
its ordinary semantic interpretation; the hardness outputs below do
not use it. The target is unrestricted validity, not truth in the
standard arithmetic structure, validity in finite structures, or
consequence from an input infinite theory.
\end{hbdefinition}

\begin{hbdefinition}{def:p04-dh-machine-encoding}{Machines used by the validity reduction}
A finite deterministic machine \(M\) has indexed finite states \(Q\),
a finite tape alphabet \(\Gamma\) containing distinct blank, zero and
one symbols, a designated initial state \(q_0\), and one halting
accepting state \(q_{\mathrm a}\). Its partial transition table consists
of entries
\[
 \delta(q,a)=(q',b,D),\qquad D\in\{\mathsf L,\mathsf R,\mathsf N\},
\]
with no entries out of \(q_{\mathrm a}\).
Missing entries do not accept. The presentation permits
\(q_0=q_{\mathrm a}\), for immediate acceptance.
States, alphabet, designated indices, table entries and the binary input
word \(w\) are explicit finite lists in the fixed string coding.
Duplicate incompatible transitions, absent designated symbols and
ill-formed words are rejected by a terminating syntactic parser.

The tape is two-way infinite and blank outside its finite written
support, as in the registered machine convention. Represent it by two
finite lists: \(\ell\) holds the symbols strictly left of the head,
nearest first; \(\rho\) holds the scanned symbol followed by those to
its right. A missing list position means blank. Lists may contain
explicit blank symbols. Let \(\mathsf a_a\) be the constructor assigned
injectively to tape symbol \(a\), and let \(\mathsf c_q\) be the binary
constructor assigned injectively to state \(q\). These are finite
selections from \(\Sigma_{\mathrm H}\), not new uninterpreted arities.
Write
\[
 [a_1\cdots a_k]
 =\mathsf a_{a_1}(\cdots\mathsf a_{a_k}(\mathsf{nil})\cdots),
 \qquad [\epsilon]=\mathsf{nil}.
\]
The initial configuration term is
\(\mathsf c_{q_0}(\mathsf{nil},[w])\).
\end{hbdefinition}

\begin{hbmethod}{def:p04-dh-horn-compilation}{The finite transition theory}
For every transition \(\delta(q,a)=(q',b,D)\), first take the scanned
nonempty-list pattern \(\rho=\mathsf a_a(v)\). Emit the universal closure
of the appropriate implication in the table below. When \(a\) is blank,
also emit the same clauses with the scanned pattern replaced by
\(\mathsf{nil}\) and its tail \(v\) replaced by \(\mathsf{nil}\).
For a left move, emit one clause for each \(d\in\Gamma\) and the separate
empty-left clause.
\[
\begin{array}{c|c|l}
D&\text{left pattern}&\text{Horn implication}\\ \hline
\mathsf R&u&
 \mathsf{Reach}(\mathsf c_q(u,\mathsf a_a(v)))
 \to\mathsf{Reach}(\mathsf c_{q'}(\mathsf a_b(u),v))\\
\mathsf N&u&
 \mathsf{Reach}(\mathsf c_q(u,\mathsf a_a(v)))
 \to\mathsf{Reach}(\mathsf c_{q'}(u,\mathsf a_b(v)))\\
\mathsf L&\mathsf a_d(u)&
 \mathsf{Reach}(\mathsf c_q(\mathsf a_d(u),\mathsf a_a(v)))
 \to\mathsf{Reach}(\mathsf c_{q'}(u,\mathsf a_d(\mathsf a_b(v))))\\
\mathsf L&\mathsf{nil}&
 \mathsf{Reach}(\mathsf c_q(\mathsf{nil},\mathsf a_a(v)))
 \to\mathsf{Reach}(\mathsf c_{q'}(\mathsf{nil},
                 \mathsf a_{\sqcup}(\mathsf a_b(v))))
\end{array}
\]
Here \(\sqcup\) is the designated blank symbol, not union. Bind exactly
the variables appearing in each clause. Let \(T_M\) be the conjunction
of all these closed clauses; an empty transition table gives
\(T_M=\topv\). Set
\[
\begin{aligned}
 I_{M,w}&=\mathsf{Reach}(\mathsf c_{q_0}(\mathsf{nil},[w])),\\
 H_M&=\exists u\,\exists v\,
             \mathsf{Reach}(\mathsf c_{q_{\mathrm a}}(u,v)),\\
 \Phi_{M,w}&=(T_M\land I_{M,w})\to H_M.
\end{aligned}
\]
Universal closure is taken before forming the conjunction and implication.
This finite sentence has no comprehension, least-fixed-point or induction
axiom. Its construction is effective; its transition-clause count is
at most \(2(|\Gamma|+1)|\delta|\), and its written size is polynomial
in the explicit machine and word encoding.
\end{hbmethod}

\begin{hblemma}{lem:p04-dh-computation-validity}{proved=annex:p04.dh.computation}{Acceptance exactly becomes unrestricted validity}
For every valid machine/word pair,
\[
 M\text{ accepts }w\quad\Longleftrightarrow\quad
 \models\Phi_{M,w}.
\]
If the machine does not accept, a countermodel is the free ground-term
algebra for \(\Sigma_{\mathrm H}\), with \(\mathsf{Reach}\) interpreted
as exactly the configuration terms in its actual run from \(w\).
Freeness and the explicit empty-list cases are part of the argument;
reachability is not merely asserted to be first-order definable.
\end{hblemma}

\begin{hbtheorem}{thm:p04-dh-validity-complete}{proved=annex:p04.dh.validity}{First-order validity is c.e.-complete}
The language \(\mathsf{VALID}_{\mathrm H}\) is computably enumerable
and c.e.-complete under total computable many-one reductions
\(\leq_m\) of \cref{sec:universal-reductions-02}. The same c.e.-completeness
therefore holds for general effectively encoded first-order validity.
For every c.e. \(A\subseteq\N\), fix a recognizer \(M_A\); an explicit
reduction is
\[
 x\longmapsto
 \operatorname{code}\bigl(\Phi_{M_A,\operatorname{bin}(x)}\bigr).
\]
It is a total syntactic construction, not a simulation run until
acceptance. The proof establishes the lower bound from every c.e. set,
rather than inferring it from mere undecidability. It also proves the
recognizer conversion needed here directly from the definition of a
c.e. domain, without promoting the sketched arithmetic-history theorem.
\end{hbtheorem}

\section{Relationalization and the finite-model boundary}
\label{sec:p04-dh-boundaries}

\begin{hbtheorem}{thm:p04-dh-relational-validity}{proved=annex:p04.dh.relational}{A purely relational validity variant}
General relational first-order validity with identity is also c.e.-complete
under total computable many-one reductions. The reduction needs only
one unary reachability predicate and finitely many binary and ternary
function-graph predicates in each input sentence; its signature is encoded
as part of the input. No statement for an arbitrary individual fixed
finite relational signature, or for the identity-free monadic fragment,
is inferred.
\end{hbtheorem}

Replace each used unary function by a binary graph predicate and each used
binary function by a ternary graph predicate. Require existence and uniqueness
of the output for every argument tuple. Replace the constant by a variable
\(z_0\), translate term values by their graph relations, and universally
quantify \(z_0\) outside
\[
 (\text{total-function graph axioms})\to
       (\text{translated }\Phi_{M,w}).
\]
The graph axioms are individually closed sentences. The value translation
and its two-way model correspondence are proved in the annex. This supplies
the effective relationalization rather than a verbal instruction alone.
Church's historical axioms-to-implication passage is the distinct source
\cite[pp.~40--41]{Church1936DecisionNote}.

\begin{hbexample}{def:p04-dh-finite-boundary}{Finite validity is a different target}
Let \(f\) be unary and \(c\) a constant, and put
\[
 B=(\forall x\,f(x)\ne c)\land
     \forall x\,\forall y\,(f(x)=f(y)\to x=y).
\]
In a finite nonempty domain, an injective self-map is surjective, so
\(B\) has no finite model. Hence \(B\to\bottom\) is true in every finite
structure. It is not unrestrictedly valid: on \(\N\), interpret
\(f(n)=n+1\) and \(c=0\). This also shows why a countermodel involving
an unbounded term domain cannot be replaced by a finite search.

Neither the c.e.-completeness proof nor this example alters
\cref{thm:trakhtenbrot}: finite validity over its stated relational
signatures is not c.e. The transition reduction uses arbitrary nonempty
structures and makes no claimed reduction to a monadic function-free
decision problem.
\end{hbexample}

\index{2-SAT!NL-completeness}\index{NL!complementation}
\index{first-order validity!computably enumerable completeness}
\index{many-one reduction!effective hardness}
